% Chunk: \S5.2, Causal Scalar Fields.
% Source: Weinberg QFT Vol. I, physical PDF pp. 224--229 (printed pp. 201--206).
% Equations: (5.2.1)--(5.2.31). Figures: none. Footnotes: 2.
% Status: source-reviewed, compile-clean, and visually checked; frozen.
% Uncertainties: none.

\subsection{Causal Scalar Fields}

We first consider one-component annihilation and creation fields
$\phi^{(+)}(x)$ and $\phi^{(-)}(x)$ that transform as the simplest of all
representations of the Lorentz group, the scalar, with
$D(\Lambda)=\mathbf{1}$. Restricted to rotations, this is just the scalar
representation of the rotation group, for which $\boldsymbol{\mathcal{J}}=0$,
so Eqs.~\eqref{eq:5.1.25} and \eqref{eq:5.1.26} have no solutions except for $j=0$, in which
case $\sigma,\bar\sigma$ take only the value zero. Thus a scalar field can only
describe particles of zero spin. Assuming also for the moment that the field
describes only a single species of particle, with no distinct antiparticle
(and dropping the species label $n$ as well as the spin label $\sigma$ and the
field label $\ell$), the quantities $u_\ell(0\sigma n)$ and $v_\ell(0\sigma n)$
are here just the numbers $u(0)$ and $v(0)$. It is conventional to adjust the
overall scales of the annihilation and creation fields so that these constants
both have the values $(2m)^{-1/2}$. Eqs.~\eqref{eq:5.1.21} and \eqref{eq:5.1.22} then give simply
\begin{equation}
  u(\mathbf{p})=(2p^0)^{-1/2}
  \tag{5.2.1}\label{eq:5.2.1}
\end{equation}
and
\begin{equation}
  v(\mathbf{p})=(2p^0)^{-1/2}.
  \tag{5.2.2}\label{eq:5.2.2}
\end{equation}
The fields \eqref{eq:5.1.17} and \eqref{eq:5.1.18} are then, in the scalar case,
\begin{equation}
  \phi^{(+)}(x)=\int \dd^3p\,(2\pi)^{-3/2}(2p^0)^{-1/2}
  a_{\mathbf{p}}e^{\ii p\cdot x}
  \tag{5.2.3}\label{eq:5.2.3}
\end{equation}
and
\begin{equation}
  \phi^{(-)}(x)=\int \dd^3p\,(2\pi)^{-3/2}(2p^0)^{-1/2}
  a^\dagger_{\mathbf{p}}e^{-\ii p\cdot x}=\phi^{(+)\dagger}(x).
  \tag{5.2.4}\label{eq:5.2.4}
\end{equation}
A Hamiltonian density $\mathcal{H}_I(x)$ that is formed as a polynomial in
$\phi^{(+)}(x)$ and $\phi^{(-)}(x)$ will automatically satisfy the requirement
\eqref{eq:5.1.9}, that it transform as a scalar. It remains to satisfy the other
condition for the Lorentz invariance of the $S$-matrix, that
$\mathcal{H}_I(x)$ commute with $\mathcal{H}_I(y)$ at space-like separations
$x-y$. If $\mathcal{H}_I(x)$ were a polynomial in $\phi^{(+)}(x)$ alone,
there would be no problem. All annihilation operators commute or anticommute,
so $\phi^{(+)}(x)$ either commutes or anticommutes with $\phi^{(+)}(y)$ for all
$x$ and $y$, according to whether the particle is a boson or fermion,
respectively:
\begin{equation}
  [\phi^{(+)}(x),\phi^{(+)}(y)]_\mp=0.
  \tag{5.2.5}\label{eq:5.2.5}
\end{equation}
Hence any $\mathcal{H}_I(x)$ formed as a polynomial in $\phi^{(+)}(x)$ (or,
for fermions, any such even polynomial) will commute with $\mathcal{H}_I(y)$
for all $x$ and $y$. The problem, of course, is that, in order to be Hermitian,
$\mathcal{H}_I(x)$ must involve $\phi^{(+)\dagger}(x)=\phi^{(-)}(x)$ as well as
$\phi^{(+)}(x)$, and $\phi^{(+)}(x)$ does not commute or anticommute with
$\phi^{(-)}(y)$ for general space-like separations. Using the commutation
(for bosons) or anticommutation (for fermions) relations \eqref{eq:4.2.5}, we have
\[
  [\phi^{(+)}(x),\phi^{(-)}(y)]_\mp
  =\int\frac{\dd^3p\,\dd^3p'}{(2\pi)^3(2p^0\cdot2p'^0)^{1/2}}
  e^{\ii p\cdot x}e^{-\ii p'\cdot y}\delta^3(\mathbf{p}-\mathbf{p}'),
\]
which collapses to the single integral
\begin{equation}
  [\phi^{(+)}(x),\phi^{(-)}(y)]_\mp=\Delta_+(x-y),
  \tag{5.2.6}\label{eq:5.2.6}
\end{equation}
where $\Delta_+$ is a standard function:
\begin{equation}
  \Delta_+(x)\equiv\frac{1}{(2\pi)^3}\int\frac{\dd^3p}{2p^0}e^{\ii p\cdot x}.
  \tag{5.2.7}\label{eq:5.2.7}
\end{equation}
This is manifestly Lorentz-invariant, and therefore for space-like $x$ it can
depend only on the invariant square $x^2>0$. We can thus evaluate $\Delta_+(x)$
for space-like $x$ by choosing the coordinate system so that
\[
  x^0=0,\qquad |\mathbf{x}|=\sqrt{x^2}.
\]
Eq.~\eqref{eq:5.2.7} then gives
\begin{align*}
  \Delta_+(x)
  & =\frac{1}{(2\pi)^3}\int\frac{\dd^3p}{2\sqrt{\mathbf{p}^2+m^2}}
  e^{\ii\mathbf{p}\cdot\mathbf{x}}
  \\
  & =\frac{4\pi}{(2\pi)^3}\int_0^\infty
  \frac{p^2\,\dd p}{2\sqrt{p^2+m^2}}
  \frac{\sin(p\sqrt{x^2})}{p\sqrt{x^2}}.
\end{align*}
Changing the variable of integration to $u=p/m$, this is
\begin{equation}
  \Delta_+(x)=\frac{m}{4\pi^2\sqrt{x^2}}
  \int_0^\infty\frac{u\,\dd u}{\sqrt{u^2+1}}
  \sin(m\sqrt{x^2}u)
  \tag{5.2.8}\label{eq:5.2.8}
\end{equation}
or, in terms of a standard Hankel function,
\begin{equation}
  \Delta_+(x)=\frac{m}{4\pi^2\sqrt{x^2}}K_1(m\sqrt{x^2}).
  \tag{5.2.9}\label{eq:5.2.9}
\end{equation}
This isn't zero, so what are we to do with it? Note that even though
$\Delta_+(x)$ is not zero, for $x^2>0$ it is even in $x^\mu$. Instead of using
only $\phi^{(+)}(x)$, suppose we try to construct $\mathcal{H}_I(x)$ out of a
linear combination
\[
  \phi(x)=\kappa\phi^{(+)}(x)+\lambda\phi^{(-)}(x).
\]
Using Eq.~\eqref{eq:5.2.6}, we have then for $x-y$ space-like
\begin{align*}
  [\phi(x),\phi^\dagger(y)]_\mp
  & =|\kappa|^2[\phi^{(+)}(x),\phi^{(-)}(y)]_\mp
  +|\lambda|^2[\phi^{(-)}(x),\phi^{(+)}(y)]_\mp
  \\
  & =(|\kappa|^2\mp|\lambda|^2)\Delta_+(x-y),
\end{align*}
\begin{align*}
  [\phi(x),\phi(y)]_\mp
  & =\kappa\lambda\bigl([\phi^{(+)}(x),\phi^{(-)}(y)]_\mp
  +[\phi^{(-)}(x),\phi^{(+)}(y)]_\mp\bigr)
  \\
  & =\kappa\lambda(1\mp1)\Delta_+(x-y).
\end{align*}
Both of these will vanish if and only if the particle is a \textit{boson}
(i.e., it is the top sign that applies) and $\kappa$ and $\lambda$ are equal
in magnitude
\[
  |\kappa|=|\lambda|.
\]
We can change the relative phase of $\kappa$ and $\lambda$ by redefining the
phases of the states so that $a_{\mathbf{p}}\to e^{\ii\alpha}a_{\mathbf{p}}$,
$a^\dagger_{\mathbf{p}}\to e^{-\ii\alpha}a^\dagger_{\mathbf{p}}$, and hence
$\kappa\to\kappa e^{\ii\alpha}$, $\lambda\to\lambda e^{-\ii\alpha}$. Taking
$\alpha=\tfrac12\operatorname{Arg}(\lambda/\kappa)$, we can in this way make
$\kappa$ and $\lambda$ equal in phase, and hence equal.

Redefining $\phi(x)$ to absorb the overall factor $\kappa=\lambda$, we have then
\begin{equation}
  \phi(x)=\phi^{(+)}(x)+\phi^{(+)\dagger}(x)=\phi^\dagger(x).
  \tag{5.2.10}\label{eq:5.2.10}
\end{equation}
The interaction density $\mathcal{H}_I(x)$ will commute with
$\mathcal{H}_I(y)$ at space-like separations $x-y$ if formed as a
normal-ordered polynomial in the self-adjoint scalar field $\phi(x)$.

Even though the choice of the relative phase of the two terms in Eq.~\eqref{eq:5.2.10}
is a matter of convention, it is a convention that once adopted must be used
wherever a scalar field for this particle appears in the interaction Hamiltonian
density. For instance, suppose that the interaction density involved not only
the field \eqref{eq:5.2.10}, but also another scalar field for the same particle
\[
  \widetilde\phi(x)=e^{\ii\alpha}\phi^{(+)}(x)
  +e^{-\ii\alpha}\phi^{(+)\dagger}(x)
\]
with $\alpha$ an arbitrary phase. This $\widetilde\phi$, like $\phi$, would be
causal in the sense that $\widetilde\phi(x)$ commutes with
$\widetilde\phi(y)$ when $x-y$ is space-like, but $\widetilde\phi(x)$ would not
commute with $\phi(y)$ at space-like separations, and therefore we cannot have
both of these fields appearing in the same theory.

If the particles that are destroyed and created by $\phi(x)$ carry some
conserved quantum number like electric charge, then $\mathcal{H}_I(x)$ will
conserve the quantum number if and only if each term in $\mathcal{H}_I(x)$
contains equal numbers of operators $a_{\mathbf{p}}$ and $a^\dagger_{\mathbf{p}}$.
But this is impossible if $\mathcal{H}_I(x)$ is formed as a polynomial in
$\phi(x)=\phi^{(+)}(x)+\phi^{(+)\dagger}(x)$. To put this another way, in
order that $\mathcal{H}_I(x)$ should commute with the charge operator $Q$
(or some other symmetry generator) it is necessary that it be formed out of
fields that have simple commutation relations with $Q$. This is true for
$\phi^{(+)}(x)$ and its adjoint, for which
\[
  [Q,\phi^{(+)}(x)]_-=-q\phi^{(+)}(x),
\]
\[
  [Q,\phi^{(+)\dagger}(x)]_-=+q\phi^{(+)\dagger}(x),
\]
but not for the self-adjoint field \eqref{eq:5.2.10}.

In order to deal with this problem, we must suppose that there are
\textit{two} spinless bosons, with the same mass $m$, but charges $+q$ and
$-q$, respectively. Let $\phi^{(+)}(x)$ and $\phi^{(+)c}(x)$ denote the
annihilation fields for these two particles, so that\footnote{The label `$c$'
denotes `charge conjugate.' It should be kept in mind that a particle that
carries no conserved quantum numbers may or may not be its own antiparticle,
with $a^c_{\mathbf{p}}=a_{\mathbf{p}}$.}
\[
  [Q,\phi^{(+)}(x)]_-=-q\phi^{(+)}(x),
\]
\[
  [Q,\phi^{(+)c}(x)]_-=+q\phi^{(+)c}(x).
\]
Define $\phi(x)$ as the linear combination
\[
  \phi(x)=\kappa\phi^{(+)}(x)+\lambda\phi^{(+)c\dagger}(x),
\]
which manifestly has the same commutator with $Q$ as $\phi^{(+)}(x)$ alone
\[
  [Q,\phi(x)]_-=-q\phi(x).
\]
The commutator or anticommutator of $\phi(x)$ with its adjoint is then, at
space-like separation
\begin{align*}
  [\phi(x),\phi^\dagger(y)]_\mp
  & =|\kappa|^2[\phi^{(+)}(x),\phi^{(+)\dagger}(y)]_\mp
  +|\lambda|^2[\phi^{(+)c\dagger}(x),\phi^{(+)c}(y)]_\mp
  \\
  & =(|\kappa|^2\mp|\lambda|^2)\Delta_+(x-y),
\end{align*}
while $\phi(x)$ and $\phi(y)$ automatically commute or anticommute with each
other for all $x$ and $y$ because $\phi^{(+)}$ and $\phi^{(+)c\dagger}$ destroy
and create different particles. In deriving this result, we have tacitly
assumed that the particle and antiparticle have the same mass, so that the
commutators or anticommutators involve the same function $\Delta_+(x-y)$.
Fermi statistics is again ruled out here, because it is not possible that
$\phi(x)$ should anticommute with $\phi^\dagger(y)$ at space-like separations
unless $\kappa=\lambda=0$, in which case the fields simply vanish. So a
spinless particle must be a boson.

For Bose statistics, in order that a complex $\phi(x)$ should commute with
$\phi^\dagger(y)$ at space-like separations, it is necessary and sufficient
that $|\kappa|^2=|\lambda|^2$, as well as for the particle and antiparticle to
have the same mass. By redefining the relative phase of states of these two
particles, we can again give $\kappa$ and $\lambda$ the same phase, in which
case $\kappa=\lambda$. This common factor can again be eliminated by a
redefinition of the field $\phi$, so that
\[
  \phi(x)=\phi^{(+)}(x)+\phi^{(+)c\dagger}(x)
\]
or in more detail
\begin{equation}
  \phi(x)=\int\frac{\dd^3p}{(2\pi)^{3/2}(2p^0)^{1/2}}
  \left[a_{\mathbf{p}}e^{\ii p\cdot x}
  +a^{c\dagger}_{\mathbf{p}}e^{-\ii p\cdot x}\right].
  \tag{5.2.11}\label{eq:5.2.11}
\end{equation}
This is the essentially unique causal scalar field. This formula can be used
both for purely neutral spinless particles that are their own antiparticles
(in which case we take $a^c_{\mathbf{p}}=a_{\mathbf{p}}$), and for particles with
distinct antiparticles (for which $a^c_{\mathbf{p}}\ne a_{\mathbf{p}}$).

For future use, we note here that the commutator of the complex scalar field
with its adjoint is
\begin{equation}
  [\phi(x),\phi^\dagger(y)]=\Delta(x-y),
  \tag{5.2.12}\label{eq:5.2.12}
\end{equation}
where
\begin{equation}
  \Delta(x-y)\equiv\Delta_+(x-y)-\Delta_+(y-x)
  \equiv\int\frac{\dd^3p}{2p^0(2\pi)^3}
  \left[e^{\ii p\cdot(x-y)}-e^{-\ii p\cdot(x-y)}\right].
  \tag{5.2.13}\label{eq:5.2.13}
\end{equation}

Let's now consider the effect of the various inversion symmetries on this
field. First, from the results of Section 4.2, we can readily see that the
effect of the space-inversion operator on the annihilation and creation
operators is:\footnote{We are omitting the subscript 0 on inversion operators
$P$, $C$, and $T$, because in virtually all cases where these inversions are
good symmetries, the same operators induce inversion transformations on `in'
and `out' states and on free-particle states.}
\begin{equation}
  P a_{\mathbf{p}}P^{-1}=\eta^*a_{-\mathbf{p}},
  \tag{5.2.14}\label{eq:5.2.14}
\end{equation}
\begin{equation}
  P a^{c\dagger}_{\mathbf{p}}P^{-1}=\eta^c a^{c\dagger}_{-\mathbf{p}},
  \tag{5.2.15}\label{eq:5.2.15}
\end{equation}
where $\eta$ and $\eta^c$ are the intrinsic parities of the particle and
antiparticle, respectively. Applying these results to the annihilation field
\eqref{eq:5.2.3} and the charge-conjugate of the creation field \eqref{eq:5.2.4}, and changing
the variable of integration from $\mathbf{p}$ to $-\mathbf{p}$, we see that
\begin{equation}
  P\phi^{(+)}(x)P^{-1}=\eta^*\phi^{(+)}(\mathcal{P}x)
  \tag{5.2.16}\label{eq:5.2.16}
\end{equation}
\begin{equation}
  P\phi^{(+)c\dagger}(x)P^{-1}
  =\eta^c\phi^{(+)c\dagger}(\mathcal{P}x),
  \tag{5.2.17}\label{eq:5.2.17}
\end{equation}
where as before $\mathcal{P}x=(-\mathbf{x},x^0)$. We see that in general
applying the space inversion to the scalar field
$\phi(x)=\phi^{(+)}(x)+\phi^{(+)c\dagger}(x)$ would give a different field
$\phi_P=\eta^*\phi^{(+)}+\eta^c\phi^{(+)c\dagger}$. Both fields are separately
causal, but if $\phi$ and $\phi_P^\dagger$ appear in the same interaction then
we are in trouble, because in general they do not commute at space-like
separations. The only way to preserve Lorentz invariance as well as parity
conservation and the hermiticity of the interaction is to require that
$\phi_P$ be proportional to $\phi$, and hence that
\begin{equation}
  \eta^c=\eta^*.
  \tag{5.2.18}\label{eq:5.2.18}
\end{equation}
That is, \textit{the intrinsic parity $\eta\eta^c$ of a state containing a
spinless particle and its antiparticle is even}. We have now simply
\begin{equation}
  P\phi(x)P^{-1}=\eta^*\phi(\mathcal{P}x).
  \tag{5.2.19}\label{eq:5.2.19}
\end{equation}
These results also apply when the spinless particle is its own antiparticle,
for which $\eta^c=\eta$, and imply that the intrinsic parity of such a particle
is real: $\eta=\pm1$.

Charge-conjugation can be handled in much the same way. From the results of
Section 4.2, we have
\begin{equation}
  C a_{\mathbf{p}}C^{-1}=\xi^*a^c_{\mathbf{p}},
  \tag{5.2.20}\label{eq:5.2.20}
\end{equation}
\begin{equation}
  C a^{c\dagger}_{\mathbf{p}}C^{-1}=\xi^c a^\dagger_{\mathbf{p}},
  \tag{5.2.21}\label{eq:5.2.21}
\end{equation}
where $\xi$ and $\xi^c$ are the phases associated with the operation of
charge-conjugation on one-particle states. It follows then that
\begin{equation}
  C\phi^{(+)}(x)C^{-1}=\xi^*\phi^{(+)c}(x),
  \tag{5.2.22}\label{eq:5.2.22}
\end{equation}
\begin{equation}
  C\phi^{(+)c\dagger}(x)C^{-1}=\xi^c\phi^{(+)\dagger}(x).
  \tag{5.2.23}\label{eq:5.2.23}
\end{equation}
In order that $C\phi(x)C^{-1}$ should be proportional to the field
$\phi^\dagger(x)$ with which it commutes at space-like separations, it is
evidently necessary that
\begin{equation}
  \xi^c=\xi^*.
  \tag{5.2.24}\label{eq:5.2.24}
\end{equation}
Just as for ordinary parity, the intrinsic charge-conjugation parity
$\xi\xi^c$ of a state consisting of a spinless particle and its antiparticle
is even. We now have simply
\begin{equation}
  C\phi(x)C^{-1}=\xi^*\phi^\dagger(x).
  \tag{5.2.25}\label{eq:5.2.25}
\end{equation}
Again, these results apply also in the case where the particle is its own
antiparticle, where $\xi^c=\xi$. In this case the charge-conjugation parity
like the ordinary parity must be real, $\xi=\pm1$.

Finally we come to time-reversal. From Section 4.2 we have
\begin{equation}
  T a_{\mathbf{p}}T^{-1}=\zeta^*a_{-\mathbf{p}},
  \tag{5.2.26}\label{eq:5.2.26}
\end{equation}
\begin{equation}
  T a^{c\dagger}_{\mathbf{p}}T^{-1}=\zeta^c a^{c\dagger}_{-\mathbf{p}}.
  \tag{5.2.27}\label{eq:5.2.27}
\end{equation}
Recalling that $T$ is antiunitary, and again changing the variable of
integration from $\mathbf{p}$ to $-\mathbf{p}$, we find that
\begin{equation}
  T\phi^{(+)}(x)T^{-1}=\zeta^*\phi^{(+)}(-\mathcal{P}x)
  \tag{5.2.28}\label{eq:5.2.28}
\end{equation}
\begin{equation}
  T\phi^{(+)c\dagger}(x)T^{-1}
  =\zeta^c\phi^{(+)c\dagger}(-\mathcal{P}x).
  \tag{5.2.29}\label{eq:5.2.29}
\end{equation}
In order for $T\phi(x)T^{-1}$ to be simply related to the field $\phi$ at the
time-reversed point $-\mathcal{P}x$, we must have
\begin{equation}
  \zeta^c=\zeta^*
  \tag{5.2.30}\label{eq:5.2.30}
\end{equation}
and then
\begin{equation}
  T\phi(x)T^{-1}=\zeta^*\phi(-\mathcal{P}x).
  \tag{5.2.31}\label{eq:5.2.31}
\end{equation}
